\section{来源审计：这不是一张万能 Alignment 配方表}

这场课最值得保留的不是若干固定超参数，而是一种实验组织方式：先把“helpful chatbot”拆成 demonstration data、preference data、training objective 与 evaluation protocol，再逐项改变数据分布和评价器，观察结论是否稳定。旧稿把这些层次压成了 SFT、reward model、PPO/DPO 的名词清单，还写入未被课堂支持的统一学习率与 batch size；本次重写回到 2023 年 10 月 31 日课堂证据，并把讲者的限定语放回结论旁边。

\subsection{官方材料与版本辨认}

本讲以 Stanford CS25 V3 课程页、Stanford Online 官方录像 \texttt{mcep6W8oB1I}、人工 \texttt{en-US} 字幕和讲者公开 deck 为主源。讲者站点存在 67 页与 71 页两个版本；录像在 00:26:30--00:26:40 明确投影了 ``Recipe 2'' 和 ``Zephyr-7B''，因此选用包含这些页面的 71 页 \texttt{transformers\_united.pdf}，而不是缺少 distillation 与 UN advisory 页的 67 页变体。

\lecturefigure{slide-01.jpg}{官方 deck 标题页：Recipes for Training Helpful Chatbots}{官方幻灯片 p.1}

\begin{knowledgebox}{读图：Recipe 指工程分解，不是神秘调参}
标题中的复数 ``Recipes'' 已经提示不存在唯一流程。读后续每组实验时都应问四个问题：样本由谁写、优化什么目标、用什么评价器、结论能否跨评价器复现。若只记录模型名称和排行榜，就会漏掉本讲真正的研究方法。
\end{knowledgebox}

\lecturefigure{slide-02.jpg}{H4 团队、目标、ingredients 与 procedure}{官方幻灯片 p.2}

\begin{importantbox}{H4 的四个 H}
H4 指 \term{Helpful、Harmless、Honest、Huggy}。本讲虽然从四项目标出发，却主要研究 helpfulness 与 SFT data；harmlessness、honesty 和完整 safety evaluation 只在 preference/red-teaming 边界中出现，不能把本讲等同于全栈安全方案。
\end{importantbox}

\teachervoice{Rajani 开场就说明：团队希望在 open-access pretrained model 上复现 alignment 的 ``secret sauce''，但课堂大部分时间聚焦第一步---怎样把基座模型变成能完成用户任务的 helpful chatbot。这个范围声明决定了后文所有结果的解释边界。}

\lecturefigure{slide-03.jpg}{课堂路线：SFT、RLHF 数据、distillation、实验与评价器偏差}{官方幻灯片 p.3}

\begin{knowledgebox}{读图：本讲有两条 recipe 与一条审计线}
第一条 recipe 是 human-curated demonstrations/preferences，第二条是 Zephyr 的 AI-generated data 与 AI feedback；两条路线最后都进入 evaluation。最右侧 GPT-4 evaluator quirks 不是附录，而是对前面所有排行榜结论的反向审计。
\end{knowledgebox}

\subsection{证据边界与覆盖策略}

71 页 deck 中有 66 页承载教学信息；纯分隔页 6、32、48、59 与最终 Thanks 页 71 有意省略，其余 progressive highlights 均保留。正文引用的榜单、成本和模型能力都是 lecture-time snapshot：例如 ``Zephyr beats ChatGPT'' 只表示它在当时显示的 AlpacaEval/MT-Bench 协议上取得特定结果，不证明更安全、更真实或在所有任务上更强。

\begin{warningbox}{不要把 2023 snapshot 写成 timeless law}
排行榜会随 prompt template、judge model、sampling、数据泄漏和参赛模型变化。本讲最强的持久结论不是某个模型排名，而是：数据实验必须和 evaluator audit 一起读；否则更长、更像 judge 训练分布的回答可能被误当成更 helpful。
\end{warningbox}

\subsection{本章小结}

本讲是一个 data--objective--evaluator 联合实验，而不是通用调参表。正确阅读方式是始终区分课堂事实、讲者工程判断、排行榜代理指标与我们基于原始论文补充的机制解释。

